\section{Memory Measurement}
\label{sec:memory}

\subsection{Scope}
\label{sec:memory-scope}

The measured quantity is post-fit reconstruction memory.  The measurement boundary starts at
integrity-checked reconstruction inputs (\cref{def:recon-inputs}) and ends when the final 3D
Gaussian file is saved.  RGB capture and the fit from RGB to 2D Gaussians are upstream of
the boundary.  They run once, offline, per view, and on hardware unconstrained by the
reconstruction GPU.  We disclose the upstream fit cost alongside the result.
\todo{Produce the upstream cost table with fit time, fit memory, and storage per view
(T-2).}  The result itself concerns the resource profile of reconstruction, which is the
stage that currently forces dataset-scale GPU memory.

\subsection{Memory model}
\label{sec:memory-model}

Let $b$ be bytes per channel, $V$ the view count, $H \times W$ the resolution, and $N$ the
3D Gaussian count.  To first order,
\begin{align}
  M_{\text{dense}} &\;\approx\; \underbrace{V \, H \, W \, 3\, b}_{\text{resident RGB working set}}
    \;+\; M_{\text{model+opt}}(N) \;+\; M_{\text{render}}(H, W),
  \label{eq:mem-dense}\\[2pt]
  M_{\text{compact}} &\;\approx\; \underbrace{\textstyle\sum_{v} N_v \, \beta}_{\text{all captures}}
    \;+\; M_{\text{model+opt}}(N) \;+\; M_{\text{query}}(A),
  \label{eq:mem-compact}
\end{align}
with $\beta \approx 40$ bytes per 2D component and $M_{\text{query}}$ the per-step sampled
batch of \cref{sec:streaming}, which is independent of $V$, $H$, and $W$.  The dense working
set grows linearly in view count and quadratically in resolution.  The compact working set
is capped by the capture budget of at most $168$~kB per view and is dwarfed by
$M_{\text{model+opt}}$.  For the development scene with $V{=}26$ views at $5328 \times 4608$
this is $7.7$~GB (float32) or $1.9$~GB (uint8) of resident RGB versus $4.4$~MB of captures.
Streaming dense images from disk instead of holding them trades $M_{\text{dense}}$ for IO
bandwidth and latency, so the protocol measures IO alongside memory.

\Cref{eq:mem-dense,eq:mem-compact} are an accounting sketch and not evidence.  Real
allocators fragment, renderers hold workspaces, and CUDA context overhead is shared.  The
result is decided exclusively by the measured protocol below.

\subsection{Protocol}
\label{sec:protocol}

Each measured cell runs in a fresh process on an otherwise idle, named GPU and records the
following quantities.
\begin{itemize}[leftmargin=1.6em,itemsep=1pt]
  \item The NVML process-used-memory peak, sampled at a frozen rate.  This is the headline
        number.  CUDA-allocator counters miss driver and backend allocations and cannot be
        the sole memory figure.
  \item PyTorch allocated and reserved peaks, reported separately for decomposition.
  \item Process RSS, bytes read and written, wall time, per-stage time, and the final
        Gaussian count.
  \item Software, driver, and GPU identity, plus the background-memory baseline before
        process start.
\end{itemize}
The rules are as follows.  One warm-up run, then at least three measured repeats per seed
and scene cell, and at least five for throughput.  Identical device, resolution,
initialization (byte-for-byte, \cref{sec:paths}), seeds, and stopping rule across arms.
Interleaved A/B execution order to cancel thermal drift.  Every raw repeat is retained,
since headline medians alone are insufficient.  The compact working-set size is reported
separately from model and optimizer memory.  The headline metrics are
\begin{enumerate}[leftmargin=1.6em,itemsep=1pt]
  \item peak reconstruction GPU memory, compact versus RGB-backed, at matched primitive and
        view counts, summarized as the reduction factor $k$,
  \item the maximum number of simultaneously held views at a fixed GPU memory budget,
  \item bytes on disk per view, and iterations per second at matched primitive count.
\end{enumerate}
The statement to establish reads as follows.  At a fixed GPU memory budget and a fixed
primitive count, the compact path holds $k$ times more views than the RGB path, at held-out
quality inside the preregistered parity band.  Parity carries the result.  A smaller but
worse reconstruction is a much weaker outcome, and the rules of \cref{sec:rules} treat it as
such.

\begin{TodoBlock}
No cell of this protocol has been run, so the entire quantitative content of this section is
pending.  Required are (i) the NVML sampler and fresh-process resource harness shared by
the compact and RGB arms (T-1), (ii) the frozen parity band and materiality rule (T-3),
(iii) at least one outcome-unseen calibrated scene, since the development frames used so
far are outcome-exposed and cannot become confirmatory by renaming (T-4), (iv) standard
benchmark scenes for external validity (T-5), and (v) the measurements themselves,
reported in \cref{fig:memory-views} and \cref{tab:main-results} (T-6).  Everything above is
protocol design, implemented mechanism, or format arithmetic.
\end{TodoBlock}

\begin{figure}[t]
  \centering
  \figslot{5.5cm}{%
    \textbf{Figure 5, the headline figure.}  Peak reconstruction GPU memory on the y-axis
    versus number of views on the x-axis at fixed primitive count.  One line for RGB-backed
    training, which grows linearly and crosses the GPU capacity line, and one line for
    compact supervision, which stays near flat.  Mark the GPU capacity as a horizontal
    line, annotate the view count where the RGB path exceeds a 24~GB GPU, and annotate the
    factor $k$ at the largest common view count.  Optionally add a second panel with
    iterations per second versus views.  Uses the protocol of \cref{sec:protocol} on one
    dome scene with view-count sweeps, NVML peaks, and repeat spread as error bars.}
  \caption{\redcaption{Peak GPU memory versus view count, dense versus compact.  This is
    the figure that carries the memory result.}}
  \label{fig:memory-views}
\end{figure}
